\section{Models and Theory}
\label{sec:models}

Phase~1 treats RUL estimation as tabular regression. Each engine cycle is a
predictor vector $\mathbf{x}_{u,t}$, either the raw cycle, settings and 21
sensors or the engineered matrix of Section~\ref{sec:feature_engineering},
and the target is the uncapped training RUL $y_{u,t}$. Three model families
are compared under one protocol: a Random Forest~\cite{Breiman2001random},
XGBoost~\cite{Chen2016xgboost} and an Elastic Net~\cite{Zou2005regularization}.
Each is wrapped in a scikit-learn pipeline whose feature step (sensor selection,
regime normalization, rolling features and, for the Elastic Net, scaling) is
refitted inside every training fold, so tuning never sees statistics from the
engines it is scored on. All three are fitted by squared error and tuned by
grid search over three engine-grouped folds with seed~42, using
cross-validated RMSE as the selection criterion. Under squared error, each
model estimates the conditional mean $\mathbb{E}[y\mid\mathbf{x}]$. The
asymmetric NASA score is reported but is not optimized, so the models are not
deliberately biased towards early predictions.

\subsection{Random Forest}
A regression tree partitions the predictor space with axis-aligned splits.
Each split is chosen greedily to minimize the squared error of the two
children, and a leaf predicts the mean target of the training rows it holds.
A deep tree has low bias but high variance: small changes in the training
engines can change the early splits and, with them, the whole tree. A Random
Forest~\cite{Breiman2001random} grows $B$ such trees, each on a bootstrap sample
of rows and with only a random subset of predictors considered at every split,
and averages their predictions. If each tree has variance $\sigma^2$ and two
trees have pairwise correlation $\rho$, the variance of the average
is~\cite{Hastie2009elements}
\begin{equation}
    \operatorname{Var}\!\left[\frac{1}{B}\sum_{b=1}^{B}T_b(\mathbf{x})\right]
    = \rho\sigma^2 + \frac{1-\rho}{B}\sigma^2 .
\end{equation}
More trees remove the second term, while random predictor subsets lower
$\rho$ and therefore the first term, which more trees cannot reduce.

The method suits C-MAPSS for three reasons. Splits capture non-linear and
interacting sensor effects without a specified functional form, such as a
sensor that only matters late in life. Trees are invariant to monotone
rescaling, so they need no standardization. Because each prediction is an
average of training targets, it stays within the observed RUL range and is
never negative. The same property means a forest cannot extrapolate beyond the
longest training life. The implementation keeps scikit-learn's unlimited depth
and bootstrap sampling, sets \texttt{max\_features}=0.5 so that each split
considers half the predictors, and tunes the number of trees
$B\in\{50,100\}$ and the minimum leaf size $\in\{1,3\}$.

\subsection{XGBoost}
Gradient boosting~\cite{Friedman2001greedy} builds an additive model
$\hat{y}^{(M)}(\mathbf{x})=\sum_{m=1}^{M}\eta f_m(\mathbf{x})$, where each new
tree $f_m$ is fitted to the negative gradient of the loss of the current
ensemble. With squared error, this gradient is the residual, so every tree
corrects what the previous trees still get wrong. Where a forest reduces
variance by averaging independent deep trees, boosting reduces bias by adding
shallow, dependent trees. XGBoost~\cite{Chen2016xgboost} fits each tree using a
second-order expansion of the loss and adds an explicit complexity penalty
\begin{equation}
    \Omega(f)=\gamma T+\tfrac{1}{2}\lambda\sum_{j=1}^{T}w_j^2 ,
\end{equation}
where $T$ is the number of leaves and $w_j$ the leaf values. For squared error
the optimal leaf value is $w_j^\ast=\sum_{i\in I_j}r_i/(|I_j|+\lambda)$, the mean
residual $r_i$ of the $|I_j|$ rows in leaf $j$ shrunk towards zero, so small
leaves move the prediction less.

XGBoost suits the problem for the same non-linearity reasons as the forest and
usually reaches lower bias, at the cost of more hyperparameters that interact.
The implementation uses the squared-error objective and the histogram tree
method. It keeps the library defaults for the learning rate ($\eta=0.3$),
$\lambda=1$, $\gamma=0$ and row and column subsampling (none), and tunes the
number of trees $M\in\{100,200\}$ and the maximum depth $\in\{3,5\}$.

\subsection{Elastic Net}
The Elastic Net~\cite{Zou2005regularization} is linear regression with a
combined penalty. In the parameterization used by scikit-learn it solves
\begin{equation}
    \min_{\beta_0,\boldsymbol{\beta}}\;
    \frac{1}{2N}\sum_{i=1}^{N}\bigl(y_i-\beta_0-\mathbf{x}_i^{\top}\boldsymbol{\beta}\bigr)^2
    +\alpha\rho\lVert\boldsymbol{\beta}\rVert_1
    +\frac{\alpha(1-\rho)}{2}\lVert\boldsymbol{\beta}\rVert_2^2 ,
\end{equation}
where $\alpha$ sets the overall strength and the ratio $\rho$ moves between
ridge ($\rho=0$) and lasso ($\rho=1$). The $\ell_1$ term sets some coefficients
exactly to zero, giving a sparse model. The $\ell_2$ term stabilizes the
solution when predictors are correlated and produces the grouping effect:
correlated predictors receive similar coefficients instead of one being
selected arbitrarily. C-MAPSS sensors that respond to the same degradation are
strongly correlated, and the engineered matrix adds three rolling summaries of
each sensor at three windows, so this property is directly relevant. Because
the penalty depends on coefficient size, predictors are standardized inside
each training fold before fitting.

The Elastic Net is the linear reference. If it matched the trees, non-linear
models would be unnecessary; its coefficients are also directly interpretable.
Its limitations follow from linearity. It cannot represent a flat RUL
relationship early in life followed by a steep decline, and its predictions
are unbounded, so a test engine can receive a negative RUL. The implementation
tunes $\alpha\in\{0.1,1.0\}$ and $\rho\in\{0.2,0.8\}$ and allows 20,000
coordinate-descent iterations.

\subsection{Bias, Variance and Regularization}
For squared error, the expected error of a fitted model $\hat{f}$ at a point
$\mathbf{x}$ splits into three parts~\cite{Geman1992neural,Hastie2009elements}:
\begin{equation}
    \mathbb{E}\bigl[(y-\hat{f}(\mathbf{x}))^2\bigr]
    =\sigma^2(\mathbf{x})
    +\bigl(\mathbb{E}[\hat{f}(\mathbf{x})]-f(\mathbf{x})\bigr)^2
    +\operatorname{Var}\bigl[\hat{f}(\mathbf{x})\bigr],
\end{equation}
the irreducible noise, the squared bias and the variance, where
$f(\mathbf{x})=\mathbb{E}[y\mid\mathbf{x}]$ and the expectation is over training
sets. Each tuned hyperparameter moves a model along this trade-off.
In the forest, a larger minimum leaf size averages more rows per leaf, raising
bias and lowering variance, and \texttt{max\_features} lowers the correlation
$\rho$. In XGBoost, greater depth and more trees lower bias and raise variance,
while the learning rate and $\lambda$ shrink each tree's contribution. In the
Elastic Net, a larger $\alpha$ shrinks all coefficients, accepting bias for
lower variance.

Table~\ref{tab:model_selection} lists the configurations selected on all 100
FD001 training engines. On engineered features, cross-validation chose the
more flexible end of each tree grid (leaf size~1; depth~5 with 200 trees),
while on raw features it chose the simpler end. The Elastic Net chose stronger
shrinkage with a larger ridge share on the engineered matrix, consistent with
its many correlated rolling features, although the fold variation is too
large to treat this as more than suggestive. The differences in CV RMSE between models and between feature sets are
smaller than the standard deviation across folds, so cross-validation alone
does not establish a preferred model.

\begin{table}[t]
\centering
\caption{Grid-search selections on the FD001 training engines (three
engine-grouped folds, seed 42). CV RMSE is the mean $\pm$ standard deviation
over folds, in cycles.}
\label{tab:model_selection}
\begin{tabular}{llll}
\hline
Model & Features & Selected & CV RMSE \\
\hline
Random Forest & raw        & $B=100$, leaf 3          & $37.6\pm3.8$ \\
Random Forest & engineered & $B=100$, leaf 1          & $36.4\pm5.0$ \\
XGBoost       & raw        & $M=100$, depth 3         & $38.1\pm3.7$ \\
XGBoost       & engineered & $M=200$, depth 5         & $38.0\pm5.1$ \\
Elastic Net   & raw        & $\alpha=0.1$, $\rho=0.8$ & $40.7\pm4.8$ \\
Elastic Net   & engineered & $\alpha=1.0$, $\rho=0.2$ & $41.8\pm3.8$ \\
\hline
\end{tabular}
\end{table}

The uncapped target adds a large noise term. FD001 training engines lived
between 128 and 362 cycles. Early in life, before degradation is visible in the
sensors, engines of the same age look alike but have very different remaining
lives, so $\sigma^2(\mathbf{x})$ is large and no model can remove it. Tree SHAP
on the engineered models attributes far more of each prediction to
\texttt{cycle} than to any sensor feature, which shows that the models fall back on
an age-based estimate where sensors carry little signal. Correspondingly, the
tree models' test RMSE is two to eight times larger for engines in the early
stage than in the late stage. This is why much of the
literature caps the training target at a constant early in
life~\cite{Heimes2008recurrent,Zheng2017long}: capping removes much of this
noise. It also changes the quantity being predicted, so this project keeps the
uncapped target as primary and treats capping as an ablation.

Two limitations of the tuning follow from this framing. Each grid has two
values per hyperparameter, and several selections lie on a grid boundary, so
the optimum may lie outside the searched range. The XGBoost learning rate,
which controls variance together with the number of trees, is fixed at its
default rather than tuned. Wider grids would let the trade-off be located
rather than bracketed.

\subsection{Machine Learning for RUL}
Tree ensembles and regularized regressors are established data-driven
prognostic tools. Wu et al.~\cite{Wu2017comparative} compared Random Forests
with feed-forward neural networks and support vector regression for tool-wear
prediction on 315 milling tests and found the forest the most accurate.
On C-MAPSS, Wang et al.~\cite{Wang2023remaining} used Random Forest importance
to select engine parameters before training a multilayer perceptron, using the
forest for feature selection rather than as the final predictor.
Mathew et al.~\cite{Mathew2017prediction} compared a range of classical
regressors on the same turbofan data, and Khelif
et al.~\cite{Khelif2017direct} used support vector regression for direct RUL
estimation. Our baselines differ from these studies by evaluating all three
families under one engine-grouped protocol with fold-fitted preprocessing,
which provides the reference for the Phase~2 sequence models.
